% !PRINTED 435
\weberpage{435}
Now, in order to determine $\omega$, we make use of the relation
\[
\omega \chi = \chi^2 \omega,
\]
and obtain, if
\[
\omega = \begin{pmatrix} \alpha_1, & \alpha_2, & \alpha_3 \\ \beta_1, & \beta_2, & \beta_3 \\ \gamma_1, & \gamma_2, & \gamma_3 \end{pmatrix}
\]
is set,
\[
\begin{pmatrix} \alpha_2, & \alpha_3, & \alpha_1 \\ \beta_2, & \beta_3, & \beta_1 \\ \gamma_2, & \gamma_3, & \gamma_1 \end{pmatrix}
=
\begin{pmatrix} \beta_1, & \beta_2, & \beta_3 \\ \gamma_1, & \gamma_2, & \gamma_3 \\ \alpha_1, & \alpha_2, & \alpha_3 \end{pmatrix},
\]
hence
\[
\begin{aligned}
\alpha_1 &= \beta_3 = \gamma_2 = \alpha, \\
\alpha_2 &= \beta_1 = \gamma_3 = \beta, \\
\alpha_3 &= \beta_2 = \gamma_1 = \gamma,
\end{aligned}
\]
and $\omega$ thus takes the form
\weq{5}{\omega = \begin{pmatrix} \alpha, & \beta, & \gamma \\ \beta, & \gamma, & \alpha \\ \gamma, & \alpha, & \beta \end{pmatrix}.}

If we express the condition that $\omega$ is to be of degree 2,
the relations
\weq{6}{\begin{aligned} \alpha^2 + \beta^2 + \gamma^2 &= 1 \\ \beta\gamma + \gamma\alpha + \alpha\beta &= 0. \end{aligned}}
result.

From (6) it follows that $(\alpha + \beta + \gamma)^2 = 1$, and if one sets the
determinant of $\omega$ equal to 1, then with the use of (6) it follows that
$(\alpha + \beta + \gamma)^3 = -1$, consequently
\weq{7}{\alpha + \beta + \gamma = -1.}

From (6) and (7) we further infer the relations:
\weq{8}{\alpha = \beta\gamma - \alpha^2, \quad \beta = \gamma\alpha - \beta^2, \quad \gamma = \alpha\beta - \gamma^2}
\weq{9}{\alpha^3 + \beta^3 + \gamma^3 - 3\alpha\beta\gamma = -1,}
of which the latter, which one obtains from
\[
(\alpha + \beta + \gamma)^3 - 3(\beta\gamma + \gamma\alpha + \alpha\beta)(\alpha + \beta + \gamma) = -1
\]
represents the negative value of the determinant of $\omega$.

Finally we further form, according to the relations [\S.\,72, (15)]
\[
\Theta = \chi \tau^3 \omega \tau^4, \quad \Theta^3 = \chi \tau^6 \omega \tau^2 \text{:}
\]
\weq{10}{\Theta = \begin{pmatrix} \gamma \varepsilon^2, & \alpha \varepsilon^6, & \beta \\ \alpha, & \beta \varepsilon^4, & \gamma \varepsilon^5 \\ \beta \varepsilon^3, & \gamma, & \alpha \varepsilon \end{pmatrix}, \quad \Theta^{-1} = \begin{pmatrix} \gamma \varepsilon^5, & \alpha, & \beta \varepsilon^4 \\ \alpha \varepsilon, & \beta \varepsilon^3, & \gamma \\ \beta, & \gamma \varepsilon^2, & \alpha \varepsilon^6 \end{pmatrix}.}
